\subsection{High-domain unit eighth moments for every single pair \((2,\allowbreak b)\)}\label{proof:Q5-C057}
Work on a nonempty finite probability space with a real antisymmetric kernel \(|K|\le\allowbreak 1\), and use
\[
\begin{gathered}
A=J+K,\ T=J-K,\ g=K1,\ p=1+g,\ q=1-g,\ d=Eg^2,\quad\\
\theta=13/32,\quad\kappa=19/32,\quad \delta_j=\|K^j1\|_2^2 .
\end{gathered}
\]
Write \(f_n=\allowbreak A^n\one\) and \(h_n=\allowbreak Tf_{n-1}\). The fixed actual direction is
\[
F_b=h_2f_b,\qquad H_b=pT F_b,\qquad b\ge4.
\]
We prove
\[
\boxed{EH_b^8\le199/200<1-d/500
\quad\text{for every integer }b\ge5,\quad1/4\le d\le1/3.} \tag{Q57.0.1}\label{eq:Q57:0.1}
\]
The finite spectral input is that of Section~\ref{proof:Q5-C081}. The finite retained-envelope construction and its high-domain estimates are given below. The low-domain bound for every \(b\ge\allowbreak 4\), and the full-domain \(b=\allowbreak 4\) bound, are proved in Section~\ref{proof:Q5-C055}. They are used only in the final union.

\paragraph{Q57.1 The retained finite construction on the full domain}
Let \(0<\allowbreak d\le\allowbreak 1/3\), \(m=\allowbreak 1-d\) and \(\tau=\allowbreak d-\delta_2\). Define
\[
\beta_1=\beta_2=0,\qquad
\beta_j=\sum_{i=1}^{\lfloor(j-1)/2\rfloor}(-1)^{i-1}\delta_i .
\]
The finite positive convolution and its alternating terms give
\(\tau\le\allowbreak \beta_j\le\allowbreak  d\) for \(j\ge\allowbreak 3\). For \(j\ge\allowbreak 5\), retain
\[
\beta_j\le d-\delta_2+\delta_3\le d-\kappa \delta_2 . \tag{Q57.1.1}\label{eq:Q57:1.1}
\]
If the finite sum ends earlier the first inequality still holds; each later pair begins with a negative term and has nonpositive sum. Put
\(R_{\mathrm{base}}=\allowbreak \sum_{j=\allowbreak 3}^bK^j\one\).
Since \(h_2=\allowbreak 2-f_2\), positive path convolution gives
\[
h_2+f_b=2+R_{\mathrm{base}}-\sum_{j=3}^b\beta_j f_{b-j}.
\]
For \(b=\allowbreak 5\), take \(C_1=\allowbreak d,\allowbreak C_2=\allowbreak \delta_2,\allowbreak C_3=\allowbreak 0\).
For \(b\ge\allowbreak 6\), take
\[
\lambda=\frac{2d-\beta_b-\beta_{b-1}}
 {\beta_{b-2}+(1-d)\beta_{b-3}},\quad
C_1=\beta_{b-1},\ C_2=\lambda\beta_{b-2},\
C_3=\lambda\beta_{b-3}.
\]
The denominator is positive and the numerator is nonnegative. The full constant from these four losses is at least
\((4-d)\tau\ge\allowbreak (11/3)\kappa d=\allowbreak 209d/96>\allowbreak 2d\), so \(\lambda\le\allowbreak 1\).
Thus the retained portion is valid throughout \(d\le\allowbreak 1/3\), and
\[
0\le C_i\le d,\quad C_1\ge\tau,\quad
\beta_b+C_1+C_2+(1-d)C_3=2d .
\]
Set
\[
C_R=C_1+C_2+C_3,\quad D_R=C_2+C_3,\quad E_R=C_3,\quad
R=R_{\mathrm{base}}-C_Rg-D_Rr-E_Rs .
\]
Here \(r=\allowbreak Kg,\allowbreak s=\allowbreak Kr\). Retaining the actual \(f_1,\allowbreak f_2,\allowbreak f_3\) terms in the positive convolution gives
\[
0\le h_2+f_b\le2m+R,\qquad
0\le F_b\le(m+R/2)^2 . \tag{Q57.1.2}\label{eq:Q57:1.2}
\]
The base before squaring is nonnegative. Moreover,
\[
C_R=2d-\beta_b+dE_R,\qquad
D_R=2d-\beta_b-\beta_{b-1}+dE_R,\qquad0\le E_R\le d .
\tag{Q57.1.3}\label{eq:Q57:1.3}
\]

\paragraph{Q57.2 The same finite spectral representation and its positive mean}
Let \(N=\allowbreak b-2\). On \(0\le\allowbreak  x\le\allowbreak \theta\), write
\[
\sum_{j=0}^{N-1}(i\sqrt x)^j=a(x)+i\sqrt x\,b(x).
\]
Finite alternating sums give \(1-x\le\allowbreak  a,\allowbreak b\le\allowbreak 1\); for \(N\ge\allowbreak 4\), also \(a,\allowbreak b\le\allowbreak 1-x+x^2\).
Let \(\nu\) be the positive spectral measure relative to \(g\), with mass \(d\), first moment \(\delta_2\) and second moment \(\delta_3\). The multiplier of \(R\) is
\[
-A(x)-i\sqrt x\,B(x),\quad
A=C_R+x(a-E_R),\quad B=D_R+xb .
\]
Since \(a\ge\allowbreak \kappa>\allowbreak d\ge\allowbreak  E_R\), both \(A,\allowbreak B\) are nonnegative.
Put \(V=\allowbreak \E R^2,\allowbreak D_0=\allowbreak \E R,\allowbreak Z=\allowbreak KR\) and
\(\alpha_{\rm actual}=\allowbreak \E Z\). Skew adjointness and the spectral representation give
\[
V=\int(A^2+xB^2)\,d\nu,\quad
D_0=\int B\,d\nu,\quad
\alpha_{\rm actual}=\int A\,d\nu . \tag{Q57.2.1}\label{eq:Q57:2.1}
\]
The mean and variance estimates below apply to this same \(R\).
By \eqref{eq:Q57:1.1} and \eqref{eq:Q57:1.3},
\(C_R\ge\allowbreak  d+\kappa\delta_2+dE_R\). Since \(a\ge\allowbreak 1-x\),
\[
\begin{aligned}
\alpha_{\rm actual}
&\ge d^2+\kappa(1+d)\delta_2-E_R(\delta_2-d^2)\\
&\ge d^2+(\kappa-\theta d)\delta_2+d^3\\
&\ge (1+\kappa+\kappa d)d^2=:\alpha .
\end{aligned} \tag{Q57.2.2}\label{eq:Q57:2.2}
\]
The second step uses \(E_R\le\allowbreak  d,\allowbreak \delta_2\ge\allowbreak  d^2\); the last uses
\(\kappa-\theta d>\allowbreak 0\).

\paragraph{Q57.3 A high-domain spectral envelope for every \(b\ge\allowbreak 5\)}
Set
\[
a_0(d)=3/8+3d/2,\qquad
z_0(d)=79/256+13d/16+d^2,\qquad U=a_0\sqrt d .
\]
For every \(b\ge\allowbreak 5\) and \(1/4\le\allowbreak  d\le\allowbreak 1/3\), we show
\[
V\le d a_0^2,\qquad D_0\le d z_0,\qquad
|Z|\le U,\qquad EZ^2\le\theta U^2 . \tag{Q57.3.1}\label{eq:Q57:3.1}
\]
The last two bounds follow from \(|KR|\le\allowbreak \|R\|_2\) and the spectral bound once the first is proved.

For \(b=\allowbreak 5\), \(C_R=\allowbreak d+\delta_2,\allowbreak D_R=\allowbreak \delta_2,\allowbreak E_R=\allowbreak 0\), and
\(A\le\allowbreak \theta-\theta^2+(1+\theta)d\), \(B\le\allowbreak \theta+\theta d\).
The triangle inequality in two coordinates yields
\[
\sqrt{A^2+xB^2}\le3/8+3d/2.
\]
The constants are
\[
(\theta-\theta^2)^2+\theta^3
=131313/1048576<9/64,
\]
\[
(1+\theta)^2+\theta^3=66997/32768<9/4 .
\]
This vector estimate holds on the full legal domain. Also
\(B\le\allowbreak \theta(1+d)\le\allowbreak  z_0\) when \(d\ge\allowbreak 1/4\): the difference is increasing, and at \(d=\allowbreak 1/4\), replacing \(79/256\) by the smaller
\(u=\allowbreak \theta-\theta^2+\theta^3\) already leaves \(2165/32768>\allowbreak 0\).

For \(b\ge\allowbreak 6\), put \(u=\allowbreak \theta-\theta^2+\theta^3=\allowbreak 10101/32768<\allowbreak 79/256\).
By \eqref{eq:Q57:1.3} and \(a,\allowbreak b\le\allowbreak 1-x+x^2\), the nonnegative upper bounds for \(A,\allowbreak B\) increase on the spectral interval: their derivatives are at least \(2/3-E_R>\allowbreak 0\) and \(2/3\). Hence
\[
A^2+xB^2\le
\{u+(1+\theta)d+E_R(d-\theta)\}^2
 +\theta\{u+2\theta d+dE_R\}^2. \tag{Q57.3.2}\label{eq:Q57:3.2}
\]
The right side is a convex quadratic in \(E_R\in[0,\allowbreak d]\).
At \(E_R=\allowbreak 0\), the triangle inequality gives
\[
\sqrt{\text{right side}}
\le u\sqrt{1+\theta}
 +d\sqrt{(1+\theta)^2+4\theta^3}
<3/8+3d/2.
\]
Here \(u<\allowbreak 5/16\), \(\sqrt{1+\theta}\le\allowbreak 6/5\), and the second squared norm is
\(18397/8192<\allowbreak 9/4\), with difference \(35/8192\).

At \(E_R=\allowbreak d\), enlarge \(u\) to \(u_*=\allowbreak 79/256\). It suffices to prove
\[
(u_*+d+d^2)^2+\theta(u_*+2\theta d+d^2)^2
\le(3/8+3d/2)^2,\quad0\le d\le1/3 . \tag{Q57.3.3}\label{eq:Q57:3.3}
\]
The right side minus the left is
\(P(d)=\allowbreak A_0+B_0d+C_0d^2-D_0^*d^3-E_0d^4\), where
\[
A_0=9/64-(1+\theta)u_*^2,\quad
B_0=9/8-2u_*(1+2\theta^2),
\]
\[
C_0=5/4-4\theta^3-2(1+\theta)u_*,
\quad D_0^*=2+4\theta^2,\quad E_0=1+\theta.
\]
Since \(d^3\le\allowbreak  d^2/3,\allowbreak d^4\le\allowbreak  d^2/9\),
\(P(d)\ge\allowbreak  L(d)=\allowbreak A_0+B_0d+(C_0-D_0^*/3-E_0/9)d^2\).
The quadratic coefficient is negative: using just
\(u_*>\allowbreak 1/4,\allowbreak D_0^*\ge\allowbreak 2,\allowbreak E_0\ge\allowbreak 1\), it is at most
\(5/4-1/2-2/3-1/9=\allowbreak -1/36\).
Thus \(L\) is concave and its minimum on the interval is at an endpoint.
The exact endpoint differences are
\[
L(0)=14067/2097152>0,
\]
\[
L(1/3)=P(1/3)=822339/169869312>0 .
\]
The second denominator is \(32\cdot2304^2\). The two squared bases are
\(1735/2304\) and \(1591/2304\); their weighted numerator is \(129233853\), whereas the target \((7/8)^2\) has numerator \(130056192\).
This proves \eqref{eq:Q57:3.3} on the entire interval.
Finally \(B\le\allowbreak  u+2\theta d+d^2\le\allowbreak  z_0\), and
\eqref{eq:Q57:2.1} gives \(D_0\le\allowbreak  dz_0\).

\paragraph{Q57.4 Positive propagation and Hermite domination}
Positive propagation and AM--GM for the two nonnegative factors give
\[
H_b\le\{m+D_0/2-Z/2+V/(4m)\}^2
\le(C-Z/2)^2,
\]
\[
C(d)=1-d+\frac d2 z_0(d)+\frac{d a_0(d)^2}{4(1-d)}. \tag{Q57.4.1}\label{eq:Q57:4.1}
\]
Since \(C\ge\allowbreak  m\ge\allowbreak 2/3\) and \(U\le\allowbreak 7/8\), \(C>\allowbreak U/2\), so the square and the increase of its base are valid.
Let \(W=\allowbreak -Z\). Then
\(\E W\le\allowbreak -\alpha,\allowbreak \E W^2\le\allowbreak \theta U^2,\allowbreak |W|\le\allowbreak  U\), and
\(\alpha^2\le\allowbreak \theta U^2\).
For \(f(w)=(C+w/2)^{16}\), use \eqref{eq:oc:negative-contact}
with \(a=\alpha\), \(v_0=\theta U^2\), and
\(b_{-\alpha}\le f(U)/(U+\alpha)^2\). The positive base follows from
\(C>U/2\), and the mean coefficient is \(f'(-\alpha)+2\alpha b_{-\alpha}\ge0\).
\[
EH_b^8\le(C-\alpha/2)^{16}
 +\mathcal R(k)(C+U/2)^{16},\quad
k=\alpha/U,\quad
\mathcal R(k)=\frac{\theta-k^2}{(1+k)^2}. \tag{Q57.4.2}\label{eq:Q57:4.2}
\]
The denominator is positive and \(\theta-k^2\ge\allowbreak 0\). On the legal range \(k\ge\allowbreak 0\),
\(\mathcal R'(k)=\allowbreak -2(k+\theta)/(1+k)^3<\allowbreak 0\).
Precisely,
\[
k(d)=\frac{d^{3/2}(51+19d)}{12(1+4d)}
\]
This increases on the high domain, since its logarithmic derivative is
\(3/(2d)+19/(51+19d)-4/(1+4d)>\allowbreak 0\).
The first and third terms already have positive difference
\((3+4d)/(2d(1+4d))\).

The other needed monotone function is \(G(d)=\allowbreak C+U/2\).
For \(d\ge\allowbreak 1/4\), \(C'\ge\allowbreak -1+z_0/2\) and \(z_0\ge\allowbreak 1/2\).
Writing \(x=\allowbreak \sqrt d\),
\((U/2)'=\allowbreak (9/8)x+(3/32)/x\ge\allowbreak 3/4\): this last function increases for \(x\ge\allowbreak 1/2\) and equals \(3/4\) at \(1/2\).
Thus \(G'\ge\allowbreak 0\), with positive terms still retained.

The function \(C\) is convex. The second derivative of \(dz_0/2\) is nonnegative, and \(da_0^2/(1-d)\) is a product of two nonnegative functions with nonnegative first and second derivatives; the remaining \(1-d\) is affine.
Hence \(C\) is bounded by the larger endpoint value on \([1/4,\allowbreak 1/3]\):
\[
C(1/4)=1779/2048<113/128,\qquad
C(1/3)=6065/6912<113/128 .
\]
Also \(\alpha\) increases and \(\alpha(1/4)=\allowbreak 223/2048\). Therefore
\[
C-\alpha/2\le3393/4096<5/6,
\qquad(C-\alpha/2)^{16}<1/16. \tag{Q57.4.3}\label{eq:Q57:4.3}
\]
The power bound follows from \((5/6)^4<\allowbreak 1/2\).
For \(d\ge\allowbreak 5/16\), \(\alpha\ge\allowbreak 22775/131072\), giving
\[
C-\alpha/2\le208649/262144<4/5,
\qquad(C-\alpha/2)^{16}<3/100 . \tag{Q57.4.4}\label{eq:Q57:4.4}
\]
For the last power bound, use \((4/5)^8<\allowbreak 17/100\) and
\(289/10000<\allowbreak 3/100\); the required cross comparison is
\(6553600<\allowbreak 6640625\).
All these bases are positive, since \(\alpha\le\allowbreak 43/216<\allowbreak 1/5\) and \(C\ge\allowbreak 2/3\).

\paragraph{Q57.5 Four rational interval payments}
Use the four closed intervals
\[
[1/4,2/7],\quad[2/7,5/16],\quad
[5/16,21/64],\quad[21/64,1/3].
\]
The endpoint estimates below apply throughout these intervals by the proved monotonicity of \(k\) and \(G\).
The exact lower endpoints give
\[
k(1/4)=223/768\ge2/7,\qquad
k(2/7)=79\sqrt{14}/882\ge1/3,
\]
\[
k(5/16)=4555\sqrt5/27648\ge11/30,\qquad
k(21/64)=693\sqrt{21}/8192\ge31/80. \tag{Q57.5.1}\label{eq:Q57:5.1}
\]
The first cross difference is \(25\). The second follows from
\(87374>\allowbreak 86436\).
For the third, use \(\sqrt5\ge\allowbreak 67/30\), since \(4500>\allowbreak 4489\), then
\(305185>\allowbreak 304128\).
For the fourth, use \(\sqrt{21}\ge\allowbreak 197/43\), since \(38829>\allowbreak 38809\), then \(682605>\allowbreak 682496\).
Thus the respective upper bounds for \(\mathcal R(k)\) are
\[
\frac{509}{2592},\quad\frac{85}{512},\quad
\frac{1957}{13448},\quad\frac{1639}{12321}. \tag{Q57.5.2}\label{eq:Q57:5.2}
\]

At \(d=\allowbreak 2/7\), \(a_0=\allowbreak 45/56\); using \(\sqrt{2/7}\le\allowbreak 15/28\),
\[
G(2/7)\le
\frac57+\frac{79}{1792}+\frac{13}{392}+\frac4{343}
 +\frac{1755}{6272}
=\frac{95097}{87808}<\frac{13}{12}. \tag{Q57.5.3}\label{eq:Q57:5.3}
\]
The final cross difference is \(340\); the square-root comparison is
\(1568<\allowbreak 1575\).
At \(d=\allowbreak 5/16\),
\[
C(5/16)=\frac{78537}{90112}
<\frac{447}{512},\qquad U/2\le\frac{243}{1024},
\]
The second bound uses \(\sqrt5\le\allowbreak 9/4\). Hence
\(G(5/16)\le\allowbreak 1137/1024<\allowbreak 10/9\), with cross difference \(7\).
The same-denominator difference for the first \(C\) comparison is \(135\).
At \(d=\allowbreak 21/64\),
\[
C(21/64)=\frac{19742627}{22544384}<\frac{897}{1024},
\qquad U/2\le\frac{2035}{8192}.
\]
These comparisons use the same-denominator difference \(5725\) and
\(\sqrt{21}\le\allowbreak 55/12\), since \(3024<\allowbreak 3025\).
Consequently \(G(21/64)\le\allowbreak 9211/8192<\allowbreak 9/8\).
At \(d=\allowbreak 1/3\), \(U/2=\allowbreak 7\sqrt3/48\); using
\(\sqrt3\le\allowbreak 26/15\), since \(675<\allowbreak 676\), gives
\[
G(1/3)\le\frac{6065}{6912}+\frac{91}{360}
=\frac{39061}{34560}<\frac{181}{160}. \tag{Q57.5.4}\label{eq:Q57:5.4}
\]
The final same-denominator difference is \(35\).
The respective sixteenth-power upper bounds for the positive nodes are
\[
4,\qquad11/2,\qquad53/8,\qquad29/4. \tag{Q57.5.5}\label{eq:Q57:5.5}
\]
Indeed,
\((13/12)^4<\allowbreak 7/5\) by \(142805<\allowbreak 145152\), so its sixteenth power is less than
\((7/5)^4=\allowbreak 2401/625<\allowbreak 4\).
Also \((10/9)^8<\allowbreak 7/3\) by \(300000000<\allowbreak 301327047\), so the sixteenth power is less than \(49/9<\allowbreak 11/2\).
Next \((9/8)^8<\allowbreak 18/7\) by \(301327047<\allowbreak 301989888\), giving
\(324/49<\allowbreak 53/8\).
Finally \((181/160)^4<\allowbreak 41/25\) by \(1073283121<\allowbreak 1074790400\), and
\((41/25)^4<\allowbreak 29/4\) follows from \(11303044<\allowbreak 11328125\).

Substitute \eqref{eq:Q57:4.3}--\eqref{eq:Q57:4.4},
\eqref{eq:Q57:5.2} and \eqref{eq:Q57:5.5} into
\eqref{eq:Q57:4.2}. The four interval bounds are
\[
\frac1{16}+4\frac{509}{2592}=\frac{1099}{1296},
\]
\[
\frac1{16}+\frac{11}{2}\frac{85}{512}=\frac{999}{1024},
\]
\[
\frac3{100}+\frac{53}{8}\frac{1957}{13448}
=1-\frac{15887}{2689600},
\]
\[
\frac3{100}+\frac{29}{4}\frac{1639}{12321}
=1-\frac{3431}{616050}.
\tag{Q57.5.6}\label{eq:Q57:5.6}
\]
Each deficit exceeds \(1/200\). For the last two, the exact comparisons are
\(3177400>\allowbreak 2689600\) and \(686200>\allowbreak 616050\); the first two are immediate.
Thus \(\E H_b^8\le\allowbreak 199/200\) for every \(b\ge\allowbreak 5\) on the whole high domain.
Since \(d/500\le\allowbreak 1/1500<\allowbreak 1/200\), this proves \eqref{eq:Q57:0.1}.

\paragraph{Q57.6 The unit catalogue and its fixed direction}
Combining the high-domain proof with Section~\ref{proof:Q5-C055}'s bound for every \(b\ge\allowbreak 4,\allowbreak d\le\allowbreak 1/4\) and its full-domain \(b=\allowbreak 4\) bound yields
\[
\boxed{EH_b^8\le1-d/500
\quad\text{for every integer }b\ge4,\text{every finite admissible kernel}.} \tag{Q57.6.1}\label{eq:Q57:6.1}
\]
Here \(H_b=\allowbreak pT(h_2f_b)\) is the same actual function in both domains.
At \(d=\allowbreak 0\), leaf messages equal one.

The long terminal arm points toward its leaf.
The intermediate vertex of the two-edge short arm points toward both the terminal and its leaf.
The unit connector points from the terminal to the receiving root, whose original unit arm points outward.
The undecorated component has \(b+4\) edges.
Genuine equal-arm pairs retain their directions \(f_\ell h_\ell\le\allowbreak 1\); positivity restores any number and multiplicity of them.
Complete reversal holds with the same original probability measure, kernel and adjoint.
